\documentclass[11pt,a4paper]{article}
\usepackage[margin=25mm]{geometry}
\usepackage{amsmath,amssymb,amsthm}
\newtheorem{lemma}{Lemma}
\newtheorem{theorem}[lemma]{Theorem}
\newcommand{\gen}[1]{\langle #1\rangle}
\title{Binary block systems and transitive subgroups}
\author{Human Compression project: a component of the L044 recovery}
\date{}
\begin{document}\maketitle

Write $G_I=\gen{\rho_j:j\notin I}$ for a group with distinguished
generators $\rho_0,\ldots,\rho_{r-1}$, and $G_i=G_{\{i\}}$.
A binary block system is a partition into two nonempty blocks which
the group permutes. The following observation uses the actual
partitions, including the action of a generator omitted from one
of their stabilizers.

\begin{lemma}\label{binary}
Suppose $G\le S_\Omega$ is primitive, where $|\Omega|=n>2$.
Let $I$ have $t$ members, suppose $G_I$ is transitive, and suppose
each $G_i$, for $i\in I$, preserves a binary block system.
Then their common refinement has exactly $2^t$ equally sized blocks.
The action of $G_I$ on them is the regular translation group
$\mathbb F_2^t$. In particular $2^t$ divides $n$.
If $n=2^t$, then $G_I$ itself is this elementary abelian group in
its regular action.
\end{lemma}
\begin{proof}
Choose a function $f_i:\Omega\to\mathbb F_2$ specifying the two
blocks of the $i$th partition. For $g\in G_i$, there is a constant
$\chi_i(g)$ such that
\[
                  f_i(gx)=f_i(x)+\chi_i(g)\quad(x\in\Omega).
\]
The restriction of each $\chi_i$ to $G_I$ is a homomorphism.
Let $F=(f_i)_{i\in I}$, and let $V$ be the image of the joint
homomorphism $\chi:G_I\to\mathbb F_2^t$. Transitivity gives
$F(\Omega)=F(x_0)+V$. It also makes all nonempty fibres of $F$
equally sized.

If $V$ were proper, some nonzero vector $(a_i)_{i\in I}$ would
annihilate it. Thus $\sum_i a_i f_i$ would be constant on $\Omega$.
Choose $s$ with $a_s=1$. For every $i\ne s$ in $I$, the generator
$\rho_s$ belongs to $G_i$, so $f_i(\rho_s x)-f_i(x)$ is constant.
The displayed linear relation then makes
$f_s(\rho_s x)-f_s(x)$ constant as well. Hence $\rho_s$ preserves
the $s$th partition. All other distinguished generators already
preserve it because they belong to $G_s$. This contradicts
primitivity of $G$. Therefore $V=\mathbb F_2^t$, proving the
refinement and divisibility assertions. If each fibre is a
singleton, the actual action is the faithful regular translation
action, proving the last assertion.
\end{proof}

\begin{lemma}\label{growth}
Let a transitive group be generated by a finite distinguished
family, and let $B$ be generated by $b$ members of that family.
If $B$ has $a$ point orbits, some transitive subgroup is generated
by those $b$ members and at most $a-1$ additional members.
\end{lemma}
\begin{proof}
As long as the current subgroup is intransitive, some distinguished
generator joins two of its orbits; otherwise every generator would
preserve every current orbit. Adjoining that generator reduces the
number of orbits by at least one. At most $a-1$ repetitions suffice.
No action of the newly adjoined generator on the old orbit set is
assumed.
\end{proof}

\begin{theorem}\label{reduction}
Let $G=A_n$ in its natural action, $n\ge16$, with a string
C-marking of rank $r\ge n/2$. Suppose that every transitive
$G_i$ preserves a binary block system. If some $G_{\{i,j\}}$
is transitive, then $n=16$ and $r=8$.

In that remaining case there is a set $J$ of five distinguished
labels such that $T=\gen{\rho_j:j\in J}$ is transitive. After a
change of coordinates its points are $(a,b)\in\mathbb F_2^3
\times\mathbb F_2$, and
\[
 T\le W=\left\{(a,b)\longmapsto(a+u,b+v(a)):
 u\in\mathbb F_2^3,\ \sum_a v(a)=0\right\}.
                                                        \tag{1}
\]
In particular $|W|=1024$. The induced marking on $T$ has the
full intersection property, with its original labels.
\end{theorem}
\begin{proof}
Put $H=G_{\{i,j\}}$. By Lemma~\ref{binary}, its two binary
partitions refine to four blocks of size $l=n/4$, and its image
on them is $\mathbb F_2^2$. Choose two distinguished generators
$x,y$ of $H$ whose images form a basis, and put $B=\gen{x,y}$.
Each $B$-orbit meets all four blocks equally, so there are at
most $l$ such orbits.

If $x,y$ commute, $B$ is a Klein four-group with faithful regular
action on the four blocks. All its point orbits have size four,
and it has exactly $l$ orbits. Lemma~\ref{growth}, applied inside
$H$, gives a transitive subgroup on at most $l+1$ distinguished
generators. Its complement has at least $r-l-1\ge l-1$ labels.
Every maximal omission at one of those labels is transitive,
since it contains that subgroup. Lemma~\ref{binary} implies
\[
                          2^{l-1}\mid 4l.               \tag{2}
\]
For $l\ge6$, $2^{l-1}>4l$: this holds at six and the ratio
increases thereafter. For $l=5$, sixteen does not divide twenty.
Thus $l=4$ is the only possibility with $n\ge16$.

If $x,y$ do not commute, at least one $B$-orbit has more than
four points. Indeed on a four-point orbit the two generators
act as their commuting translations on the four blocks. Hence
$B$ has at most $l-1$ orbits. Lemma~\ref{growth} now supplies
a transitive subgroup on at most $l$ distinguished generators,
containing $B$. Its omitted label set has size at least $l$.
For $l\ge5$ this contradicts $2^l>4l$ and Lemma~\ref{binary}.
For $l=4$, that omitted set must have exactly four members;
the last assertion of Lemma~\ref{binary} makes the transitive
subgroup abelian, contradicting the presence of $x,y$.
This excludes the noncommuting case altogether.

We therefore have $n=16$ and a transitive subgroup on a set
$J_0$ of at most five distinguished labels. Put $t=r-|J_0|$.
Lemma~\ref{binary} gives $t\le4$, so $r\le9$. If $r=9$,
necessarily $|J_0|=5$ and $t=4$. The subgroup is then
elementary abelian of order sixteen. The intersection property
makes its five distinguished generators independent, impossible
in a vector space of dimension four. Thus $r=8$.
Enlarge $J_0$, if necessary, to a set $J$ of five labels;
its generated subgroup remains transitive.

For the three omitted labels, Lemma~\ref{binary} gives eight
pairs, indexed by their three binary coordinates $a$. The
subgroup $T$ acts on these pairs by translations $a\mapsto a+u$.
Choose the two point names $b$ in each pair arbitrarily. Every
permutation inducing a translation has the form in (1).
A translation on the eight pairs is even on the sixteen points:
each exchanged pair of pairs contributes two transpositions.
The remaining sign is $(-1)^{\sum_a v(a)}$. Since $T\le A_{16}$,
the stated even-sum condition holds. There are eight translations
and $2^7$ permitted functions $v$, proving the order assertion.
Intersection properties on subsets of $J$ are inherited from
the original C-marking.
\end{proof}

\paragraph{The remaining finite obligation.}
Degree sixteen still requires completions of transitive five-mark
subgroups in $W$. The $\binom85=56$ label sets $J$ give twenty-eight
cases up to reversal. Identify the omitted labels with the three
coordinates of $a$. Each omitted generator preserves the other
two binary partitions, so belongs to the even part of
$(S_4)^4\rtimes\mathbb F_2^2$ on their four cells. It must break
its own partition. String relations, generation and all subset
intersections remain obligations; no exhaustion is asserted here.

\paragraph{Attribution and dependencies.}
The target is Cameron--Fernandes--Leemans--Mixer,
\emph{Highest rank of a polytope for $A_n$}, arXiv:1605.09173v1,
Proposition~3.4. The source uses intersecting binary partitions
and a four-partition divisibility argument. Lemma~\ref{binary}
supplies their general independence proof; orbit growth replaces
the separate large-degree and degree-twenty/twenty-four arguments.
Other imprimitive branches must supply our binary-system hypothesis.
The case where every double omission is intransitive remains separate.
No priority for binary partitions, finite computation, conjecture,
Lean verification or whole-original credit is asserted.
\end{document}
